\documentclass[crop=false]{standalone}
\usepackage{15150toolbox}
\usepackage{comment}
\usepackage{needspace}

% Toggle: set \soltrue to render solutions inline (blue boxes).
\newif\ifsol
\soltrue   % comment this line for the student version

\ifsol
  \newenvironment{solution}
    {\par\color{solution_color}\begin{framed}\textbf{Solution.}\par}
    {\end{framed}}
\else
  \excludecomment{solution}
\fi

% Lecture-slide macros that aren't defined in 15150toolbox.
\newcommand{\term}[1]{\textbf{\textit{#1}}}
\newcommand{\thmBox}[2]{\par\noindent\begin{mdframed}[backgroundcolor=blue!8,hidealllines=true]\textbf{Theorem#1.} #2\end{mdframed}}
\newcommand{\bcBox}[1]{\par\noindent\textbf{Base case.} #1\par}
\newcommand{\ihBox}[2]{\par\noindent\textbf{Inductive hypothesis#1.} #2\par}
\newcommand{\isBox}[2]{\par\noindent\textbf{Inductive step#1.} #2\par}

% The title sits in a top-aligned parbox so long titles wrap within the
% left column instead of colliding with the right-aligned time/points.
\newcommand{\quizproblem}[3]{%
  \needspace{5\baselineskip}%
  \stepcounter{problem}%
  \ifsol
    \subsection*{\parbox[t]{0.70\linewidth}{\raggedright Problem \arabic{problem}. #1}\hfill \timing{#2} \quad (#3 pts)}%
  \else
    \subsection*{\parbox[t]{0.82\linewidth}{\raggedright Problem \arabic{problem}. #1}\hfill (#3 pts)}%
  \fi
}

\newcommand{\subpts}[1]{\textbf{(#1 pts)}}

\newcommand{\answerspace}[1]{\ifsol\else\vspace{#1}\par\fi}

\printout{Quiz 3}

\begin{document}

\begin{center}
  {\Large \textbf{15-150 --- Quiz 3}} \\
  \vspace{4pt}
  Topics: \textit{Polymorphism}; \textit{Higher-Order Functions}; \textit{Combinators and Staging} \\
  \vspace{4pt}
  \textit{60 minutes.}
\end{center}

\vspace{1em}

\noindent\textbf{Name:} \hrulefill \hspace{2em} \textbf{Andrew ID:} \hrulefill

\vspace{1em}

\noindent For reference, throughout this quiz:

\begin{codeblock}
  map    : ('a -> 'b) -> 'a list -> 'b list
  filter : ('a -> bool) -> 'a list -> 'a list
  foldl  : ('a * 'b -> 'b) -> 'b -> 'a list -> 'b
  foldr  : ('a * 'b -> 'b) -> 'b -> 'a list -> 'b

  Int.toString : int -> string
  String.size  : string -> int
\end{codeblock}

% =============================================================================
% PROBLEM 1 --- Manual type inference
% =============================================================================

\quizproblem{Polymorphic Playtime}{5}{4}

For each declaration below, determine the \textbf{most general type} of
the function being declared.
Showing your work is optional.

\vspace{0.5em}

\textbf{(a)} \subpts{1}

\begin{codeblock}
  fun swap (x, y) = (y, x)
\end{codeblock}

\answerspace{2cm}

\begin{solution}
\code{swap : 'a * 'b -> 'b * 'a}

\textbf{Rubric:} all-or-nothing, renaming OK. \code{'a * 'a -> 'a * 'a}
(over-unified): 0.
\end{solution}

\vspace{0.5em}

\textbf{(b)} \subpts{3}

\begin{codeblock}
  fun applyAll [] x = x
    | applyAll (f :: fs) x = applyAll fs (f x)
\end{codeblock}

\answerspace{3cm}

\begin{solution}
\code{applyAll : ('a -> 'a) list -> 'a -> 'a}

The base clause returns \code{x}, so the result type equals the type of
\code{x} --- call it \code{'a}. In the recursive clause, \code{f x} is
passed back into the \code{x} position, so \code{f x : 'a} as well,
forcing \code{f : 'a -> 'a}.

\textbf{Rubric:} 2 pts \code{('a -> 'a) list} domain (the unification
is the point --- \code{('a -> 'b) list}: 0 here); 1 pt the
\code{-> 'a -> 'a} remainder, curried.
\end{solution}

% \vspace{0.5em}

% \needspace{12\baselineskip}
% \textbf{(c)} \subpts{6} This one computes nothing of value, but it does
% have a type:

% \begin{codeblock}
%   fun quux (f, x, y, k) =
%     case f x of
%       [] => (x, k ())
%     | z :: _ => (y, z)
% \end{codeblock}

% \answerspace{4cm}

% \begin{solution}
% \code{quux : ('a -> 'b list) * 'a * 'a * (unit -> 'b) -> 'a * 'b}

% Step by step: \code{f x} is cased against list patterns, so
% \code{f : 'a -> 'b list} (with \code{x : 'a}, \code{z : 'b}). The two
% branches must agree: \code{(x, k ())} $\sim$ \code{(y, z)} forces
% \code{y : 'a} and \code{k : unit -> 'b}. Result type \code{'a * 'b}.

% \textbf{Rubric:} 2 pts \code{f : 'a -> 'b list}; 2 pts the cross-branch
% unifications (\code{y} with \code{x}, \code{k : unit -> 'b}); 2 pts
% assembling the 4-tuple domain and \code{'a * 'b} result. Renaming OK.
% \end{solution}

% =============================================================================
% PROBLEM 2 --- Types of partial applications
% =============================================================================

\clearpage
\quizproblem{Curry On}{13}{20}

\textbf{(a)} \subpts{8} Write the type of each expression below. (All
of them are well-typed.)

\vspace{0.5em}

\begin{center}
\begin{tabular}{|l|p{6.5cm}|}
  \hline
  \textbf{Expression} & \textbf{Type} \\
  \hline
  \code{map Int.toString}                          & \\[2em] \hline
  \code{foldl (fn (x, acc) => acc + String.size x) 0} & \\[2em] \hline
  \code{foldr (fn (x, acc) => acc)} & \\[2em] \hline
  \code{Int.toString o (fn (x, y) => y)} & \\[2em] \hline
\end{tabular}
\end{center}

\begin{solution}
\begin{itemize}
  \item \code{map Int.toString} : \code{int list -> string list}
  \item \code{foldl (fn (x, acc) => acc + String.size x) 0} : \code{string list -> int}
  \item \code{foldr (fn (x, acc) => acc)} : \code{'a -> 'b list -> 'a}
  \item \code{Int.toString o (fn (x, y) => y)} : \code{'a * int -> string}
\end{itemize}

\textbf{Rubric:} 2/row, all-or-nothing. Renamed type variables OK;
extra or missing arrows: 0.
\end{solution}

\vspace{0.5em}

\textbf{(b)} \subpts{4} For each declaration below, write the
\textbf{most general type} of the function being declared.

\vspace{0.5em}

\begin{center}
\begin{tabular}{|l|p{6.5cm}|}
  \hline
  \textbf{Declaration} & \textbf{Type} \\
  \hline
  \code{fun cons x xs = x :: xs}                       & \\[2em] \hline
  \code{fun mystery (x, y) = fn z => if z then x else y} & \\[2em] \hline
\end{tabular}
\end{center}

\begin{solution}
\begin{itemize}
  \item \code{cons} : \code{'a -> 'a list -> 'a list}
  \item \code{mystery} : \code{'a * 'a -> bool -> 'a}
\end{itemize}

For \code{mystery}: the argument is a \emph{tuple}, the result is a
curried function of \code{z : bool}, and both branches of the
\code{if} force \code{x} and \code{y} to share a type.

\textbf{Rubric:} 2/row, all-or-nothing, renaming OK. Common misses:
\code{cons} as \code{'a * 'a list -> 'a list} (tupled --- it's
curried): 0; \code{mystery} as \code{'a * 'b -> bool -> 'a}
(missed the unification): 0.
\end{solution}

\vspace{0.5em}

\needspace{14\baselineskip}
\textbf{(c)} \subpts{6} Consider the following specification:

\spec
  {pair}
  {'a -> 'b -> 'a * 'b}
  {\code{true}}
  {\code{pair x y} $\eeq$ \code{(x, y)}}

Write \textbf{three different declarations} of \code{pair} --- two
using \code{fun} and one using \code{val}. (All three should be
equal up to syntactic sugar.)

\answerspace{10cm}

\begin{solution}
\begin{codeblock}
  fun pair x y = (x, y)

  fun pair x = fn y => (x, y)

  val pair = fn x => fn y => (x, y)
\end{codeblock}

The first is syntactic sugar for the second, which is what the third
binds directly.

\textbf{Rubric:} 2/form. The three must be genuinely different shapes
(sugared \code{fun}, \code{fun} returning a lambda, \code{val} bound to
nested lambdas) --- the same shape twice with renamed variables earns
credit once.
\end{solution}

\vspace{0.5em}

\needspace{14\baselineskip}
\textbf{(d)} \subpts{2} True or false: all functions of
type \code{t1 -> t2 -> t3} for some types \code{t1}, \code{t2}, and \code{t3} are
total, since they are curried.

\answerspace{2cm}

\begin{solution}
\textbf{False.} Syntactically-sugared curried functions
(\code{fun f x y = ...}) are trivially total, but not every value of a
curried \emph{type} is --- e.g.\ from lecture:

\begin{codeblock}
  fun f x = let val _ = loop () in fn y => 0 end
\end{codeblock}

which loops forever before returning the inner lambda.

\textbf{Rubric:} 1 pt ``false''; 1 pt a witness or the
sugar-vs-arbitrary-value distinction. Bare ``false'': 1 of 2.
\end{solution}


% =============================================================================
% PROBLEM 2 --- Fold value questions
% =============================================================================

\clearpage
\quizproblem{Fold, Right There}{3}{4}

Recall the two folds from lecture:

\begin{codeblock}
  fun foldl f z [] = z
    | foldl f z (x::xs) = foldl f (f (x, z)) xs

  fun foldr f z [] = z
    | foldr f z (x::xs) = f (x, foldr f z xs)
\end{codeblock}

\textbf{(a)} \subpts{2} What does \code{foldl (op^) "" ["1", "5", "0"]}
evaluate to? (Recall \code{^} concatenates strings.)

\answerspace{2cm}

\begin{solution}
\code{"051"} --- \code{foldl} combines with the first element first, so
the accumulator builds up in reverse:
\code{"0" ^ ("5" ^ ("1" ^ ""))}.

\textbf{Rubric:} all-or-nothing. \code{"150"}: 0.
\end{solution}

\vspace{0.5em}

\textbf{(b)} \subpts{2} What does \code{foldr (op^) "" ["1", "5", "0"]}
evaluate to?

\answerspace{2cm}

\begin{solution}
\code{"150"} --- \code{foldr} is the ``natural'' fold: the elements stay
in order, joined by the combining function:
\code{"1" ^ ("5" ^ ("0" ^ ""))}.

\textbf{Rubric:} all-or-nothing.
\end{solution}

\vspace{0.5em}

% \textbf{(c)} \subpts{4} Write a function \code{stutter} satisfying:

% \spec
%   {stutter}
%   {'a list -> 'a list}
%   {\code{true}}
%   {\code{stutter [x1, x2, ..., xn]} $\eeq$
%    \code{[x1, x1, x2, x2, ..., xn, xn]}}

% \begin{constraint}
%   Your function \textbf{must not} be recursive --- use a fold.
% \end{constraint}

% \answerspace{5cm}

% \begin{solution}
% \begin{codeblock}
%   fun stutter L = foldr (fn (x, acc) => x :: x :: acc) [] L
% \end{codeblock}

% (\code{foldr} keeps the elements in order; \code{foldl} would produce
% the \emph{reversed} stutter --- exactly the (a)/(b) distinction.)

% \textbf{Rubric:} 2 pts combine fn (\code{x :: x :: acc}); 1 pt base
% \code{[]}; 1 pt picked \code{foldr}. Used \code{foldl} (reversed):
% cap 2. Recursive: cap 1.
% \end{solution}

\newpage

% =============================================================================
% PROBLEM 3 --- HOF production: count
% =============================================================================

\vspace{1em}
\quizproblem{The One Where I Have You Implement A Bunch Of Functions Using HOFs}{12}{20}

\textbf{(a)} \subpts{4} Write a function \code{count} satisfying:

\spec
  {count}
  {('a -> bool) -> 'a list -> int}
  {\code{p} is total}
  {\code{count p L} evaluates to the number of elements \code{x} in
   \code{L} such that \code{p x} $\eeq$ \code{true}}

\begin{constraint}
Your function \textbf{must not} be recursive.
\end{constraint}

Fun fact, there's at least three different ways to write this function
using HOFs. I wonder if you can find all of them?

\answerspace{5cm}

\begin{solution}
Any of:

\begin{codeblock}
  fun count p = foldr (fn (x, acc) => if p x then 1 + acc else acc) 0
  fun count p = foldl (fn (x, acc) => if p x then acc + 1 else acc) 0
  fun count p = length o filter p
\end{codeblock}

Only one solution is required; the ``fun fact'' is flavor, not a task.

\textbf{Rubric:} 3 pts counting logic (tests \code{p}, counts exactly
the hits); 1 pt HOF machinery (base value, argument order, currying).
Recursive: cap 1.
\end{solution}

\textbf{(b)} \subpts{4} Write a function \code{double} satisfying:

\spec
  {double}
  {int list -> int list}
  {\code{true}}
  {\code{double L} evaluates to a list where each element is doubled}

\begin{constraint}
Your function \textbf{must not} be recursive.
\end{constraint}

\answerspace{4cm}

\begin{solution}
\begin{codeblock}
  val double = map (fn x => 2 * x)
\end{codeblock}

Eta-expanded forms (\code{fun double L = map (fn x => 2 * x) L}) and
\code{fn x => x + x} are equally good.

\textbf{Rubric:} 3 pts correct doubling via \code{map} (or any
equivalent HOF); 1 pt honors the constraint (non-recursive). Reversed
result: cap 1.
\end{solution}

\newpage

\textbf{(c)} \subpts{6} Write a function \code{hasNegativeBalance} satisfying:

\spec
  {hasNegativeBalance}
  {int -> int list -> bool}
  {\code{true}}
  {Given a starting balance and a list of positive-or-negative
  transactions \code{L} (applied in order), \code{hasNegativeBalance
  balance L} evaluates to \code{true} iff at some point the running
  balance becomes negative}

\begin{constraint}
  Your function \textbf{must not} be recursive.
\end{constraint}

\answerspace{5cm}

\begin{solution}
The accumulator must carry the \emph{running balance}, not just a
boolean. Two clean ways:

\begin{codeblock}
  (* freeze the balance once it goes negative *)
  fun hasNegativeBalance balance L =
    foldl (fn (x, b) => if b < 0 then b else b + x) balance L < 0

  (* or carry both the balance and the flag *)
  fun hasNegativeBalance balance L =
    let
      val (_, neg) =
        foldl (fn (x, (b, neg)) => (b + x, neg orelse b + x < 0))
              (balance, balance < 0) L
    in
      neg
    end
\end{codeblock}

\textbf{Rubric:} 3 pts accumulator tracks the running balance
(\code{bool}-only accumulator: cap 2); 2 pts detects \emph{any}
intermediate negative (final balance only: cap 3); 1 pt fold
args/initial value right. Either reading of the starting balance OK.
\end{solution}

\textbf{(d)} \subpts{6} Write a function \code{tlistmax} satisfying:

\spec
  {tlistmax}
  {int list -> int option}
  {\code{true}}
  {Given a list of integers \code{L}, \code{tlistmax L} either evaluates to
  \code{NONE} if the list is empty, or \code{SOME x} where \code{x} is the maximum element of \code{L}}

\begin{constraint}
  Your function \textbf{must not} be recursive.
\end{constraint}

\answerspace{5cm}

\begin{solution}
Fold with an \code{option} accumulator:

\begin{codeblock}
  fun tlistmax L =
    foldl (fn (x, NONE) => SOME x
            | (x, SOME m) => SOME (if x > m then x else m))
          NONE L
\end{codeblock}

Also fine: pattern-match the empty case first (two clauses, still
non-recursive), seeding the fold with the head:

\begin{codeblock}
  fun tlistmax [] = NONE
    | tlistmax (x :: xs) =
        SOME (foldl (fn (y, m) => if y > m then y else m) x xs)
\end{codeblock}

\textbf{Rubric:} 4 pts max via fold (combine function compares and
keeps the larger); 2 pts empty $\to$ \code{NONE}. Seed of \code{0}
(wrong on all-negative lists): 0 on the max piece.
\end{solution}



% =============================================================================
% PROBLEM 3 --- Staging fill-in
% =============================================================================

\clearpage
\quizproblem{Stage Directions}{9}{10}

Recall \code{msort} from lecture, which sorts a list ascending
according to a comparison function, with $O(n \log n)$ work:

\begin{codeblock}
  msort : ('a * 'a -> order) -> 'a list -> 'a list
\end{codeblock}

Also consider \code{nth}, which retrieves the element at index \code{i}
(zero-indexed), with $O(\code{i})$ work:

\begin{codeblock}
  fun nth ([] : 'a list, i : int) : 'a option = NONE
    | nth (x :: xs, 0) = SOME x
    | nth (_ :: xs, i) = nth (xs, i - 1)
\end{codeblock}

A student combines them to answer \emph{``what is the \code{k}-th
smallest element?''} queries:

\begin{codeblock}
  fun kthSmallest (L : int list) (k : int) : int option =
    let
      val sorted = msort Int.compare L
    in
      nth (sorted, k)
    end
\end{codeblock}

\textbf{(a)} \subpts{3} The student runs:

\begin{codeblock}
  val kth = kthSmallest scores
  val a = kth 0
  val b = kth 1
  val c = kth 2
\end{codeblock}

How many times is \code{msort} fully evaluated (i.e., how many times is
\code{scores} sorted) by these four declarations, in total? Explain in
\textbf{one sentence}.

\answerspace{2cm}

\begin{solution}
\textbf{3 times.} \code{kthSmallest scores} immediately evaluates to a
lambda expression without entering the \code{let}, so the sort re-runs
every time \code{kth} is applied to an index.

\textbf{Rubric:} 1 pt ``3''; 2 pts reason (\code{let} body runs only
once both curried args are supplied). ``1'': 0 --- that's the
misconception under test.
\end{solution}

\vspace{0.5em}

\needspace{14\baselineskip}
\textbf{(b)} \subpts{5} Write the \emph{staged} version of \code{kthSmallest},
in which the list is sorted as soon as \code{L} is supplied:

\answerspace{10cm}

\begin{solution}
\begin{codeblock}
  fun kthSmallest (L : int list) : int -> int option =
    let
      val sorted = msort Int.compare L
    in
      fn k => nth (sorted, k)
    end
\end{codeblock}

\textbf{Rubric:} 2 pts \code{msort Int.compare L} bound in a \code{let}
\emph{before} the lambda; 2 pts returns \code{fn k => nth (sorted, k)}
(must be a lambda --- a second curried argument on the \code{fun} line
puts the sort back inside and defeats the staging: cap 1 overall, since
that's the unstaged original); 1 pt structure/types correct.
\end{solution}

\vspace{0.5em}

\textbf{(c)} \subpts{2} Using the staged version from part (b): how
many times is \code{scores} sorted by the same four declarations from
part (a)?

\answerspace{2cm}

\begin{solution}
\textbf{Once} --- when \code{kthSmallest scores} is evaluated to define
\code{kth}, the list is sorted once.

\textbf{Rubric:} all-or-nothing.
\end{solution}

% =============================================================================
% PROBLEM 4 --- Spot the bug: trivial totality
% =============================================================================

% \vspace{1em}
% \quizproblem{A Total Misunderstanding}{5}{6}

% Recall the definition of \code{map}:

% \begin{codeblock}
%   fun map f [] = []
%     | map f (x::xs) = f x :: map f xs
% \end{codeblock}

% A student writes the following argument:

% \begin{quote}
%   \textbf{Claim:} \code{map (fn x => x div 0) [1, 2, 3]} evaluates to a
%   value.
%   \begin{enumerate}
%     \item \code{map} is total: for every function value
%       \code{f : t1 -> t2}, the application \code{map f} immediately
%       evaluates to a lambda expression, which is a value.
%     \item In particular, \code{map (fn x => x div 0)} evaluates to a
%       value of type \code{int list -> int list}.
%     \item Since \code{map} is total and \code{[1, 2, 3]} is a value,
%       \code{map (fn x => x div 0) [1, 2, 3]} evaluates to a value.
%   \end{enumerate}
% \end{quote}

% \textbf{(a)} \subpts{2} Exactly one numbered step makes an invalid
% inference. Which one?

% \answerspace{1cm}

% \textbf{(b)} \subpts{2} In \textbf{one sentence}, explain why that step
% is wrong.

% \answerspace{2cm}

% \textbf{(c)} \subpts{2} What actually happens when
% \code{map (fn x => x div 0) [1, 2, 3]} is evaluated?

% \answerspace{1.5cm}

% \begin{solution}
% \textbf{(a)} Step 3.

% \textbf{(b)} The totality of \code{map} only says \code{map f} is a
% value for every \code{f}; it says nothing about the totality of
% \code{map f}, which is a \emph{different function} --- and here
% \code{map f} is not total, because \code{f} itself is not.

% \textbf{(c)} It raises the exception \code{Div}: \code{map} applies
% \code{fn x => x div 0} to the element \code{1}, and \code{1 div 0}
% raises \code{Div}.

% \textbf{Rubric:}
% \begin{itemize}
%   \item 2 pts: identifies step 3 (all-or-nothing)
%   \item 2 pts: explanation distinguishes totality of \code{map} from
%     totality of \code{map f}. Merely restating ``\code{f} isn't total''
%     without connecting it to why step 3's inference fails: 1 of 2.
%   \item 2 pts: raises \code{Div} (1 of 2 for ``it raises an exception''
%     or ``it loops/fails'' without naming \code{Div})
% \end{itemize}
% \end{solution}

\end{document}
